\section*{الفصل \ref{ch:g1:materials-around-us} --- مِمَّ تُصنع الأشياء؟}

\begin{solution}{exo:g1:materials-around-us:1}
الكرسي \omterm{def:g1:materials-around-us:object}{جسم}. والخشب \omterm{def:g1:materials-around-us:material}{مادة}: يمكن أن يُصنع منه الكرسي.
\end{solution}

\begin{solution}{exo:g1:materials-around-us:2}
لوح النافذة: زجاج. المسمار: معدن. الحصاة: حجر. الكتاب: ورق. الجورب:
قماش (صوف أو قطن).
\end{solution}

\begin{solution}{exo:g1:materials-around-us:3}
الخشب: قلم رصاص، وباب (أو كرسي، أو طاولة). البلاستيك: صحن، ودلو
(أو لعبة، أو قارورة).
\end{solution}

\begin{solution}{exo:g1:materials-around-us:4}
الشوكة المعدنية: فهي الوحيدة التي ليست من الخشب ولا من البلاستيك.
(اثنان من الخشب وواحدة من البلاستيك.)
\end{solution}

\begin{solution}{exo:g1:materials-around-us:5}
الزجاج: فهو شفاف.
\end{solution}

\begin{solution}{exo:g1:materials-around-us:6}
المسمار الفولاذي. فالمغناطيس يجذب الحديد والفولاذ، لا الزجاج ولا الخشب.
\end{solution}

\begin{solution}{exo:g1:materials-around-us:7}
جسمان يطفوان (قطعة الخشب والغطاء البلاستيكي)؛ وثلاثة تغوص
(الحصاة وكرة الزجاج والمسمار).
\end{solution}

\begin{solution}{exo:g1:materials-around-us:8}
الشفرتان من المعدن، والمقبضان من البلاستيك. فالمقص مصنوع من
مادتين.
\end{solution}

\begin{solution}{exo:g1:materials-around-us:9}
في الحلقة الزرقاء: فهي مصنوعة من المعدن. وليست في الحلقة البرتقالية
لأن المغناطيس لا يجذبها: فهي مصنوعة من الألومنيوم، لا من الحديد ولا
من الفولاذ.
\end{solution}

\begin{solution}{exo:g1:materials-around-us:10}
يجب أن تُدخل النافذة الضوء، وأن نرى من خلالها. والزجاج شفاف؛ أما
الخشب فلا.
\end{solution}

\begin{solution}{exo:g1:materials-around-us:11}
«نعم»: المسمار ومشبك الورق والمفتاح، أي 3 أجسام. «لا»:
الحصاة والملعقة والعلبة وكرة الزجاج والجورب، أي 5 أجسام.
$3 + 5 = 8$: لكل \omterm{def:g1:materials-around-us:object}{جسم} مكانه.
\end{solution}
